% !TEX root = ../../../main.tex
% Appendix entry: detection.cwt_ridges -- "CWT ridges" (Appendix A).
\entryfield{Algorithm}
\code{_run} in \code{Adapters/detection_cwt_ridges.py} calls \code{detect} in
\code{Working/Detection/cwt_ridges.py}, the project's own implementation of the ridge-line peak
detector of \citet{du2006improved}. It does \emph{not} use \code{scipy.signal.find_peaks_cwt},
which is not the paper's algorithm (its noise is the 10th percentile of the signed smallest-scale
coefficients and its SNR numerator the smallest-scale coefficient;
\code{docs/research/analysis-blocks-primary-sources.md}, \S7), nor scipy's \code{cwt} or
\code{ricker}, which scipy~1.17 removed. The Ricker wavelet and the transform are written in the
module, with FFT convolution from \code{scipy.fft}; scipy's \code{find_peaks},
\code{maximum_filter1d} and \code{median_filter} are used for maxima and medians.

``Spikes and troughs at every scale'' means: the transform is computed at $S$ scales, and at
each scale the local maxima of the coefficients (spikes) and of the negated coefficients
(troughs) are both traced into ridges. An event of any duration within the scale range can
then be found, each with its own scale.

\begin{enumerate}
\item \emph{Scales.} $S$ = \emph{n\_scales} geometric scales from \emph{min\_scale\_s} to
  \emph{max\_scale\_s}, in samples $a_j = s_j f_s$. Scales with $a_j > n/4$ are dropped and
  counted; fewer than three left is refused by the adapter.
\item \emph{Wavelet} (\code{ricker}): the Mexican hat at scale $a$, normalised by $1/a$ (the
  paper uses $1/\sqrt a$; departure A0), truncated to $[-5a, 5a]$, discrete mean removed.
  Under $1/a$ an event's peak coefficient is about $0.84\times$ its amplitude whatever its width.
\item \emph{Edges} (A4): NaNs are first filled by linear interpolation (A11); the span is then
  extended by odd (point) reflection, hinged on the median of the first and last
  $\lceil a_1\rceil$ samples, so a straight drift continues past the ends.
\item \emph{Transform} in chunks of $2^{21}$ samples with overlap enough for the largest wavelet,
  ridge wander, baseline and noise windows (A3), streamed one scale at a time. Each row has its
  running median over $\pm 10a$ subtracted (A12), which removes curved drift the hat does not
  cancel.
\item \emph{Maxima} at each scale: columns that are the maximum of their $\pm a$ window and
  positive; the same on the negated row for troughs.
\item \emph{Ridges} (\code{_trace_ridges}), from the largest scale down: each ridge takes the
  nearest maximum at the next smaller scale within $0.5a$; a contested maximum goes to the
  nearer ridge; a ridge ends after \emph{max\_gap\_scales} consecutive misses; an unclaimed
  maximum starts a new ridge.
\item \emph{Events} (\code{_candidates}): a ridge needs at least \emph{min\_ridge\_scales}
  points. Every interior maximum of its coefficient profile over scale is a candidate (A6, so a
  fast event on a slow one gives two); a maximum at either end of the ridge is counted, not
  reported. A candidate must reach \emph{min\_snr} times the noise and stand at least one noise
  above the saddle to its neighbour.
\item \emph{Noise} (A15): $1.96\times1.4826\times$ the median absolute deviation of the
  coefficients \emph{at the candidate's own best scale}, in a window of
  $\max(\texttt{noise\_window\_s}, \pm30a)$ around it less $\pm4a$, floored at the span's
  smallest-scale noise scaled by $\sqrt{a_1/a}$. Two passes: pass~1 accepts at a fixed SNR of 3;
  pass~2 re-measures with pass~1's events left out of the windows and judges at \emph{min\_snr}.
  The paper's noise, the 95th percentile of $|$coefficients$|$ at the smallest scale, is kept per
  event as the diagnostic \code{noise_min_scale}.
\item \emph{Side lobes, steps, duplicates}: a candidate is dropped as a side lobe of an
  opposite-polarity event if, at its best scale within 2.5 scales, the opposite side reaches
  $1/0.6$ of its coefficient (A7); flagged \code{step_pair} if the opposite side is between 0.6
  and $1/0.6$ of it (A14); of two same-polarity events closer than half a scale with best scales
  within $1.5\times$, the stronger is kept (A8).
\item \emph{Location and width}: the location is the ridge's position at the smallest scale of
  the event's band whose coefficient is still at least half the peak (A9); the best scale is
  refined by a parabola in log-scale (A10); $\texttt{duration\_s} = 2\sqrt{\ln2}\,\times
  \texttt{best\_scale\_s} \approx 1.665\times$, the FWHM of a Gaussian bump whose response peaks
  at that scale. The span is $\texttt{peak} \pm \texttt{duration\_s}$, clipped to the span.
\end{enumerate}
Each departure from the paper (A0--A15) is set out, with its reason and measurements, in the
docstring of \code{Working/Detection/cwt_ridges.py}.

\entryfield{Assumptions}
\begin{itemize}
\item \emph{Units.} Not unit-aware, and needs not be: \emph{min\_snr} is a ratio of
  coefficients to noise measured on the same trace. \code{peak_coef} and \code{noise} are in the
  stored unit.
\item \emph{Sampling rate.} Scales are in seconds; \emph{min\_scale\_s} must be at least one
  sample. The tests and calibration are at 1\,Hz.
\item \emph{Baseline.} A constant and a straight drift cancel exactly; curved drift is removed
  per scale by the running median. No detrend block is needed.
\item \emph{Missing values.} NaNs are bridged by linear interpolation for the transform; an
  event whose support ($\pm5a$) touches a filled sample is flagged \code{nan_in_support}; the
  count is \code{meta["n_nan"]}. A span more than half NaN, or all NaN, is refused in words.
\item \emph{Span length.} The largest usable scale is $n/4$ samples.
\end{itemize}

\entryfield{Outputs}
A \texttt{SpanSet}, one span per event, label \texttt{spike} or \texttt{drop}, score the SNR.
Its feature table holds \code{peak_idx} (span-relative), \code{best_scale_s},
\code{duration_s}, \code{cwt_duration_s} (the same, under a name \blockref{interrogation.event_shape}
does not overwrite), \code{ridge_scales}, \code{snr}, \code{peak_coef}, \code{noise},
\code{noise_min_scale}, \code{polarity} ($\pm1$), \code{ridge_polarity}, \code{edge_flag},
\code{step_pair}, \code{nan_in_support}. \code{meta} holds the counts (by polarity, flagged,
NaN, chunks, every rejection reason), the scales used and dropped, the rules in words, the
ridges of up to 500 strongest events as (time, scale) points, a coarse scalogram (at most 1024
columns per scale, as integers times \code{unit}), and, when nothing is found, an
\code{absence} sentence giving the reasons.

\entryfield{Cost}
Linear in the span at a fixed scale range: each chunk costs $O(\text{chunk}\log\text{chunk}
\times S)$, memory is $O(\text{chunk})$, never $O(nS)$. The docstring measures 19--23\,s for
2.6 million samples at the defaults on the development machine, about 48\,s under load; the
estimate (cost model exponent 1) is empty until calibrated. No span cap.

\entryfield{Failure modes}
\begin{itemize}
\item \emph{Composites} (A13): at scales near the spacing of events, the hat spans several and
  a ridge can describe their arrangement as one slow event; most reach an edge and are flagged,
  not all. Keep \emph{max\_scale\_s} below the usual spacing.
\item \emph{Edge ridges}: on drift and noise alone, 0--4 edge-flagged events a day (mean 1.0)
  over 20 seeds; none unflagged.
\item \emph{Dense trains}: neighbours' opposite lobes pull the best scale down, so durations
  shrink (480\,s events every 600\,s read about 181\,s), and when events cover more than half of
  every noise window at their scale they become the noise: 24\,s drops every 120\,s, only 9--13
  of 718 found.
\item \emph{Slow events on wandering noise} sit at low SNR at their own scale and are lost
  (random walk: 4.7 of 7 amplitude-1 events found).
\item \emph{Lopsided events}: the centre sits towards the slow side (A5); onset needs
  \blockref{interrogation.event_shape}.
\item \emph{A genuine opposite event} close to and much weaker than its neighbour is dropped as
  a side lobe.
\end{itemize}

\entryfield{Validation}
\code{tests/test_cwt_ridges.py} checks answers, not only plumbing, on synthetic 1\,Hz days:
every event of three widths (24\,s, 8\,min, 43\,min), both polarities, and a sharkfin is found
on a curved drift ten times the smallest event, with the right polarity and duration within
$1.5\times$ (three seeds); every extra event is flagged; drift and noise alone give no event
away from the edges; a fast event on a slow one gives two events; side lobes are not reported;
a train of 24\,s drops every 300\,s is fully recalled; chunked and unchunked runs agree; the
duration rule is calibrated on isolated Gaussian bumps. On real recordings, the
\texttt{cwt\_events} template (this block for drops, then \blockref{interrogation.event_shape})
beat a 20-draw phase-randomised null on three 24\,h spans: 24, 38 and 6 events against null
means of 0.2--0.3; on 4--6\,h spans it found 0--2, too few to tell
(\code{docs/prompts/blocks/reports/W1-evaluation.md}). No test compares its events with
hand-marked ones.

\entryfield{References}
\citet{du2006improved}. The paper's steps and scipy's differences:
\code{docs/research/analysis-blocks-primary-sources.md}, \S7. Build report:
\code{docs/prompts/blocks/reports/W1-C-cwt-ridges.md}.

\entryfield{Tests}
\code{tests/test_cwt_ridges.py} (the core, above); \code{tests/test_adapter_cwt_ridges.py}
(one span per event with its features, span-relative indices, meta contents and size, the
absence sentence, polarity, refusal of a too-short span); \code{tests/test_template_cwt_events.py}
(the template).
